% ===================================================================
% Paper 1, v16 (2026-07) [+ cross-review fixes]
% cross-review (ChatGPT): abstract wording fixed -- 'dark-energy-scale
% mechanism terminated by direct lattice measurement' (was ambiguous);
% v16.4: topological-DM route (negative quartic + omega repulsion) WITHDRAWN
%   as a live direction -- companion gate P5-OMEGA-01 ran and FAILED
%   (cutoff-scale collapse, no continuum window at any coupling).
% future directions (v16.3): new topological-DM route (negative quartic +
% endogenous omega repulsion) registered as the live direction with
% open gates.
% v16: (a) DE narrative synced to companion v6.1 post-C2a: the a0
% coincidence stated as between two observed numbers, owing nothing
% to the framework; monopole sector's DE-generator candidacy noted as
% terminated by direct measurement; joint falsifiability recorded as
% cashed in the double-negative direction (abstract, sec:a0,
% conclusions). (b) rc lower-boundary censoring check executed
% (rc_floor.py, all 175): 16/26 of the [2,5) bin migrate below 2 kpc
% on an extended grid (8 to the new wall at 0.13 kpc, preferring no
% scale at all); 10th pct 2.0 -> 0.16 kpc, span [0.13, 734] kpc;
% non-universality strengthened; new robustness subsection; table
% footnote and fig. 2 caption updated; censoring caveat attached to
% small-rc statements.
% v15: data gates G1+G2 executed (repo branch gates-g1-g2, scripts
% a0_profile_likelihood.py / heldout_validation.py). (1) profiled a0
% with continuous nuisance polish: 0.95 [0.920,0.979]e-10; grid-MAP
% staircase noise (~+-10 units vs Delta=1) identified as the defect
% of the withdrawn slice -- which it INFLATED (v14's 'biased narrow'
% direction corrected); (2) formal stat interval (+-0.03) vs 5-fold
% empirical spread (+-0.10) discrepancy reported as galaxy-level
% heterogeneity; (3) held-out 5-fold: RAR beats refitted Yukawa
% out of sample by Delta(-2lnL) = -4207 despite kernel test-time
% freedom; NFW comparison framed by the freedom asymmetry; split-half
% instability reported as a sample-size requirement; (4) abstract
% numbers synced; completed gates removed from future list.
% v14: MAGI round (ChatGPT) + executed gates. (1) L0 -> schematic
% channel-dependent action, no complete-basis universal G claim;
% (2) sec:micro framed as the hypothesis being tested and falsified;
% (3) parity/mixing two-sided: parity-preserving -> exactly zero,
% optimistic admixture -> 1e-122; (4) C2d executed and reported:
% monopole gradient response = exact sine-Gordon staircase -- correct
% crossover scale, wrong exponent (X^1/2 vs X^3/2), no source; the
% emergent-MOND reading fails at leading order; a0 relation labelled
% an unexplained structural coincidence; (5) M-gates recorded:
% monotone environment triggers closed (5 kills); non-monotonic
% baryon-window (shell/torus) response registered untested (W1-W5);
% (6) m_sigma phrased as tied to m_dyn (>> kpc^-1 at the operating
% point), not unconditionally 'cutoff scale'; (7) halo subsection
% retitled as squared-profile ansatz; (8) abstract rewritten to the
% four-step arc and trimmed; (9) data gates registered: fully
% profiled a0, held-out prediction, outer-cutoff discriminant.
% v13: (a) theory narrative synced to the mediator audit and
% topological funnel of the companion dark-energy paper (v5.8):
% parity/no-go recorded at eq. muc-origin (eps ~ 1e-122 now
% load-bearing), P(X) pairing route reported as executed-and-closed
% (two-body theorem), K(X) address moved to the monopole-gas gradient
% response (pre-registered, kill criterion); (b) a0 stat error bar
% +0.12/-0.01 withdrawn -- the Delta(-2lnL)=1 slice is conditioned on
% best-fit f_int and grid-quantized (both bias narrow); systematic
% band (0.96-1.28) quoted as dominant; (c) ac-test arithmetic
% corrected: rc-Vmax slope is sample-unstable (0.51/0.82/1.39), the
% cherry-picked 0.82->1.18 line removed, instability itself reported
% as a failure mode; robust statement: a_c tracks Vmax; (d) hard
% fixes: duplicate bibitem, stale mu in [0,10], undefined s, NFW
% slope description (equal parameter count), cH0/2pi attribution
% with full-M_Pl convention, (G M_B / R_B c^2) dimension, orphan
% configuration-selection paragraph headed, roadmap completed,
% AIC rounding harmonized, order-of-magnitude wording corrected.
% v12: single-a0 refit performed and reported (Sec. refit): RAR-type
% acceleration-dependent kernel, per-galaxy mu/D/i nuisances (Table 1),
% global intrinsic-scatter likelihood (no error floor). a0 = 0.96e-10
% (0.96-1.28 across error models); Yukawa decisively inferior (+4859);
% NFW retains a moderate edge (-1614), honest dual-statement reported.
% Repo: scripts rar_refit.py, nuisance_comparison.py,
% intrinsic_scatter_comparison.py (all validation-gated).
% v11 (reframe): M/L-prior robustness completed -- the finite-r_c
% population collapses under a physical stellar normalization
% (58% -> 7% lognormal / 2% mu=1 fixed); free-mu results relabeled
% as kernel-shape diagnostics; dAIC definition stated (same-config);
% mu range corrected (mu >= 0, unbounded); microscopic-status
% paragraph added (companion m_sigma measurement places the scalar
% at the cutoff scale -> angular Goldstone / K(X) route); title
% updated accordingly. Companions: Paper 2 v2.14, Paper 3 v3.10,
% Paper 4 v5.5, Paper 5 v0.19.
% ===================================================================
\documentclass[11pt,a4paper]{article}

\usepackage{amsmath,amssymb,amsfonts}
\usepackage{geometry}
\usepackage{physics}
\usepackage{slashed}
\usepackage{hyperref}
\usepackage{booktabs}
\usepackage[numbers]{natbib}
\usepackage{graphicx}
\usepackage{longtable}
\usepackage{xcolor}

\geometry{margin=1in}

\hypersetup{
  colorlinks=true,
  linkcolor=blue,
  citecolor=blue,
  urlcolor=blue
}

\title{\textbf{A Condensate-Motivated Yukawa Halo Kernel Confronts \\
SPARC Rotation Curves: Prior Fragility, an Emergent \\
Acceleration Scale, and a Single-$a_0$ Refit}}

\author{Zeta Hoi-Ho Cheng \\
\textit{Independent Researcher, Ontario, Canada}}

\date{\today}

\begin{document}

\maketitle

\begin{abstract}
We subject a microscopically motivated dark-halo hypothesis---a
fixed-range Yukawa kernel, $V_{\rm halo}^2 = B(1-e^{-2mr})/r$, from
a lattice fermion condensate---to rotation-curve tests across all
175 SPARC galaxies.
Free fits appear acceptable ($\chi^2/\mathrm{pt} = 1.08$), but
three failure modes converge on one diagnosis: the
per-galaxy range spans 1.4 decades and trends with velocity; the
characteristic acceleration $a_c = V^2(r_c)/r_c$ tracks $V_{\max}$
(slope $1.09\pm0.17$) instead of clustering at a constant; and under
physical mass-to-light priors the finite-$r_c$ population collapses
from 58\% to 7\%.  The kernel's functional form fails: fixed length where the data
demand acceleration dependence.
The refit this diagnosis demands---an acceleration-dependent
kernel, single global $a_0$, nuisance priors, and a global
intrinsic-scatter likelihood---yields
$a_0 = 0.95\pm0.03\,(\mathrm{stat})\times10^{-10}$~m~s$^{-2}$,
error-model systematics dominating (up to $1.28$): one global
parameter beats the $2\times175$-parameter Yukawa kernel by
$\Delta\mathrm{AIC} = +4859$ in sample and by $-4207$ on held-out
galaxies, where the kernel refits per galaxy and $a_0$ is frozen.
NFW retains a moderate total-AIC advantage ($-1614$), an honest
residual.
The scale relation
$a_0 \sim \Lambda_{\rm DE}^2/M_{\rm Pl} \sim cH_0/2\pi$ relates two
\emph{observed} numbers and owes nothing to the framework; the
framework's own mechanisms for both sides are computed and closed:
the galactic response at leading order (pairing, loops, and the
monopole gradient response---right crossover scale, wrong exponent,
no source), and the dark-energy-scale mechanism terminated by direct lattice
measurement in the companion paper.  Joint falsifiability was cashed in the double-negative direction;
the coincidence outlives both mechanisms, and the falsification
chain, single-$a_0$ measurement, and NFW residual are the data-side
record.
\end{abstract}


% ===================================================================
\section{Introduction}
\label{sec:intro}
% ===================================================================

The rotation curves of spiral galaxies provide one of the most direct
probes of the distribution of matter on galactic scales.
The Spitzer Photometry and Accurate Rotation Curves (SPARC) database
\cite{Lelli:2016zqa} compiles high-quality rotation curves for 175
disk galaxies with near-infrared photometry, enabling systematic
studies of the mass discrepancy problem and the radial acceleration
relation \cite{Lelli:2017vgz,McGaugh:2016leg}.

Within the cold dark matter (CDM) paradigm, the standard approach
models dark matter halos using empirically motivated density profiles.
The Navarro--Frenk--White (NFW) profile \cite{Navarro:1996gj},
derived from $N$-body simulations, remains the default baseline,
though it faces well-known challenges including the cusp--core
problem in low-mass galaxies \cite{deBlok:2009sp,Bullock:2017xww}
and the diversity problem \cite{Oman:2015xda}.
Comprehensive halo model comparisons using the SPARC sample have
been performed by \cite{Li:2020iib}.

Several theoretical programmes seek to derive gravitational and
dark-sector phenomena from underlying microscopic degrees of freedom.
The induced gravity programme, initiated by \cite{Sakharov:1967pk}
and developed extensively in the condensed matter analogy by
\cite{Volovik:2003fe}, demonstrates that gravitational dynamics can
emerge as collective effects of quantum fields.
Analogue gravity models \cite{Barcelo:2005fc} provide concrete
laboratory realizations of emergent spacetime concepts.
Verlinde's entropic gravity \cite{Verlinde:2016toy} offers an
alternative thermodynamic derivation.
Within particle physics, Nambu--Jona-Lasinio (NJL) models
\cite{Nambu:1961tp,Nambu:1961fr} and their extensions
\cite{Klevansky:1992qe} provide a well-studied framework for
dynamical symmetry breaking and collective boson emergence via
Hubbard--Stratonovich (HS) linearization
\cite{Hubbard:1959ub,Stratonovich:1957}.

Yukawa-type corrections to the Newtonian gravitational potential
have been studied extensively as phenomenological probes of modified
gravity.
\cite{deAlmeida:2018kwq} fitted a Yukawa fifth force to 40 SPARC
galaxies, finding improved $\chi^2$ and a best-fit range
$\lambda = 5.61 \pm 0.91$~kpc.
\cite{Henrichs:2021bnd} tested Yukawa gravity using the Milky Way
rotation curve with Bayesian model selection.
More recently, \cite{Hassan:2025yuk} extended the analysis to
include all galactic components with Yukawa corrections.
In all these studies, the Yukawa potential is introduced as an
ad hoc phenomenological modification to gravity---a ``fifth force''
with free coupling and range parameters.

The present work differs fundamentally in that the Yukawa halo
profile is not postulated but \emph{derived} from a specific
microscopic construction: a lattice fermion model with four-fermion
interactions, where the scalar Hubbard--Stratonovich channel
naturally produces a massive scalar field with Yukawa-type coupling
to matter.
The halo amplitude $B$ and stellar mass-to-light rescaling $\mu$
remain phenomenological parameters at the galactic level, but the
\emph{functional form} of the halo profile has a structural origin.

Fuzzy dark matter (FDM) models
\cite{Hu:2000ke,Hui:2016ltb,Schive:2014dra}
employ ultralight scalar fields ($m \sim 10^{-22}$~eV) to produce
halo cores through quantum pressure.
Modified gravity approaches such as MOND
\cite{Milgrom:1983ca,Famaey:2011kh}
and its relativistic extensions \cite{Bekenstein:2004ne}
address the mass discrepancy through modifications of gravitational
dynamics.

This paper focuses exclusively on the scalar dark matter channel
and its galactic phenomenology.
Questions of emergent gravity (tensor channel), dark energy (vector
channel), and gauge interactions are deferred to future work.

The paper is organized as follows.
Section~\ref{sec:micro} derives the scalar halo profile from the
microscopic framework.
Section~\ref{sec:method} describes the fitting methodology.
Section~\ref{sec:results} presents results.
Section~\ref{sec:nfw} compares with NFW.
Section~\ref{sec:robust} discusses robustness.
Section~\ref{sec:discuss} provides discussion.
Section~\ref{sec:refit} performs the single-acceleration-scale
refit that the diagnosis demands, and
Section~\ref{sec:conclusions} concludes.


% ===================================================================
\section{Microscopic Framework and Emergent Scalar Profile}
\label{sec:micro}
% ===================================================================

\emph{Status of this section.}  We formulate here, in full, the
fixed-range scalar realization that motivated the original kernel.
The lattice and symmetry audits reported later in this paper and in
the companion papers show that this realization cannot provide the
required galactic mediator; it is retained as the microscopic
hypothesis being tested---and, by the end of this paper,
falsified---rather than as a construction the programme still
defends.

\subsection{Lattice fermion substrate}

We consider a discrete lattice of fermionic degrees of freedom at a
characteristic spacing $a \simeq \ell_{\mathrm{Pl}}$, governed by a
four-fermion interaction of the NJL type:
\begin{equation}
\mathcal{L}_0 = \bar{\psi}(i\slashed{\partial})\psi
+ \sum_A \frac{G_A}{2N}(\bar{\psi}\Gamma_A\psi)^2,
\label{eq:L0}
\end{equation}
where $\Gamma_A$ labels Dirac--flavor channels with
channel-dependent couplings $G_A$ and $N$ is the number of fermion
species.  The symmetry-selected condensate sector actually adopted
is specified in the companion papers; no equality among
Dirac-channel couplings is assumed here, and the present paper uses
only the scalar projection.

The vacuum is characterized by a fermion condensate that gaps the
microscopic substrate at the lattice scale, so that no
single-fermion excitations are observable at low energies; the
physical low-energy degrees of freedom are collective condensate
modes.
The light mode relevant for the present work is the phase (angular)
fluctuation of the complex condensate, whose small mass is
protected by an approximate $U(1)$ symmetry broken only by the
discrete lattice structure.
The dynamical analysis of the vacuum and the gravitational sector
is developed in a separate publication.
For the galactic-scale phenomenology developed in this paper,
the only input required from the microscopic sector is the
hypothesized existence of a massive scalar collective mode with a Yukawa-type
coupling to matter; the results below do not depend on the
details of the vacuum construction.

\subsection{Hubbard--Stratonovich scalar channel}

The scalar component of the four-fermion interaction is linearized
by introducing an auxiliary field $\chi$:
\begin{equation}
\mathcal{L}_{\mathrm{HS}}^{(S)} =
\bar{\psi}(i\slashed{\partial})\psi
+ g\chi\,\bar{\psi}\psi
- \frac{\chi^2}{2G_S},
\label{eq:HS}
\end{equation}
where $G_S$ is the scalar-channel coupling and $g$ is the Yukawa
coupling.

Integrating out the fermion fields at one loop in the large-$N$
limit generates the effective potential:
\begin{equation}
V_{\mathrm{eff}}(\chi) = \frac{\chi^2}{2G_S}
- N \int \frac{d^4 p}{(2\pi)^4}\,
\ln\!\left(p^2 + g^2\chi^2\right).
\label{eq:Veff}
\end{equation}

The condensate $\bar{\chi}$ is determined by the stationarity
condition $\partial V_{\mathrm{eff}}/\partial\chi = 0$, which
yields the gap equation
\begin{equation}
\frac{1}{G_S} = 2Ng^2 \int \frac{d^4 p}{(2\pi)^4}\,
\frac{1}{p^2 + g^2\bar{\chi}^2},
\label{eq:gap}
\end{equation}
with the (Euclidean) momentum integral regulated at the lattice
cutoff $\Lambda \sim 1/a$.
The curvature of the effective potential at the minimum,
$\partial^2 V_{\mathrm{eff}}/\partial\chi^2\big|_{\chi=\bar{\chi}}$,
sets the mass of the heavy radial (amplitude) mode.
The light mode relevant for halo phenomenology is the angular
(phase) fluctuation of the complex condensate, whose mass
$m_\chi$ is protected by the approximate phase symmetry of the
substrate and is treated as an effective parameter in this work.

\subsection{Nonlinear scalar equation}

The condensate $\chi$ satisfies the static field equation
\begin{equation}
-\nabla^2 \chi + m_\chi^2\,\chi + 4\lambda\,\chi^3 = \mu_c\,J_B,
\label{eq:scalar-field}
\end{equation}
where $\lambda$ parameterizes the quartic self-interaction and
$J_B$ is the baryonic source density; here $\chi$ denotes the light
(angular) collective mode, with mass $m_\chi$.  The monopole
coupling $\mu_c$ to baryons is treated as an effective parameter,
but its microscopic origin can be made explicit.  A pure phase
(Goldstone) mode couples only \emph{derivatively}, through
$(\partial_\mu\theta/f)\,J^\mu$, and is therefore not sourced by a
static baryon density.  The monopole coupling arises precisely
because the chiral $U(1)$ is \emph{explicitly broken}: the same
subleading breaking $\varepsilon$ that lifts the flat direction and
gives the pseudo-Goldstone its mass, $m_\chi^2 \sim \varepsilon
\Lambda^2$, also mixes the angular mode with the heavy radial
amplitude, which \emph{does} carry the tree-level monopole vertex
$g\chi\bar\psi\psi$.  Seen through this small mixing angle
$\sim\varepsilon$, the angular mode inherits a suppressed
non-derivative coupling,
\begin{equation}
\mu_c \;\sim\; \varepsilon\,g \;\sim\; \frac{m_\chi^2}{\Lambda^2}\,g ,
\label{eq:muc-origin}
\end{equation}
so that the halo source strength and the field mass share a single
origin in the explicit breaking; $\mu_c$ is not an independent vertex
but the radial monopole coupling viewed through the
breaking-induced mixing.  A first-principles evaluation of
$\varepsilon$ on the lattice substrate is developed in the companion
paper on the gravitational sector; for the galactic-scale fits below,
$\mu_c$ enters only through the effective halo normalization $B$.
Eq.~\eqref{eq:muc-origin} carries a falsifiable consequence---a
correlation between the halo amplitude and the field mass---that a
future joint analysis of the two could test.
We record its status with both branches closed.  In a
parity-preserving realization the scalar--pseudoscalar mixing---and
with it the linear coupling---vanishes identically.  Even allowing
the most optimistic parity-violating admixture, its natural size is
bounded by $\varepsilon \sim m_\chi^2/\Lambda^2 \sim 10^{-122}$,
irrelevant for galactic dynamics.  This two-sided failure is a
fourth, independent argument against the literal fixed-range
reading---anticipating the verdict of the fits below---and is
developed into a general no-go for every loop-level coupling route
in the companion dark-energy paper.

\subsection{Linearized solution}

In the linearized regime ($\lambda \to 0$) with a spherically
symmetric, localized source $J_B(r)$, the static Green's function
of Eq.~\eqref{eq:scalar-field} gives the Yukawa-type solution:
\begin{equation}
\chi(r) = \frac{\mu_c}{4\pi}\int d^3r'\,
\frac{e^{-m_\chi|\mathbf{r}-\mathbf{r}'|}}{|\mathbf{r}-\mathbf{r}'|}
\,J_B(r').
\label{eq:yukawa-sol}
\end{equation}

For a point source $J_B = M_B\,\delta^3(\mathbf{r})$, this reduces to
\begin{equation}
\chi(r) = \frac{\mu_c M_B}{4\pi}\,\frac{e^{-m_\chi r}}{r}.
\label{eq:yukawa-point}
\end{equation}

The point-source solution motivates the radial form of the halo
kernel used in the fitting analysis below.
A fully baryon-resolved derivation using the spatially extended
source distribution $J_B(r)$ from the observed gas and stellar
profiles is left for future work.
Such a calculation would relate the halo amplitude directly to
the baryonic mass, potentially eliminating $B$ as a free parameter.

The Yukawa length scale $r_c = 1/m_\chi$ is motivated by the
scalar mass arising from the effective
potential~\eqref{eq:Veff} and gap equation~\eqref{eq:gap}.
In a complete microscopic theory, $m_\chi$ would be a universal
constant determined by the lattice parameters.
In the present phenomenological analysis, $m$ is treated as an
effective galactic-scale parameter: it is microscopically motivated
but not yet predicted from first principles.
Whether the data support a single universal $m$ or require
galaxy-dependent values is tested quantitatively in
Section~\ref{sec:universal-m}.

\subsection{Squared-profile halo ansatz and circular-velocity kernel}

The scalar condensate carries an energy density whose leading
contribution in the mass-dominated regime is
$\rho_\chi \propto m_\chi^2 |\chi|^2$.
This approximation neglects the gradient term
$(\nabla\chi)^2$, which can be comparable to $m^2\chi^2$
at small radii.
Equation~\eqref{eq:rho-halo} below should therefore be understood
as a leading scaling ansatz---a phenomenological kernel motivated
by the squared Yukawa profile, not a derivation from the full
$T^{00}$;
a derivation from the full energy-momentum tensor $T^{00}$
of the scalar field, with all terms retained, is deferred to
future work.
We note that the singular small-radius behavior of both the
field profile and its energy density is an artifact of the
point-source idealization: for a spatially extended baryonic
source, the convolution in Eq.~\eqref{eq:yukawa-sol} regularizes
the profile at radii below the source scale, which provides
further motivation for the extended-source treatment.
For the Yukawa solution~\eqref{eq:yukawa-point}, this gives
\begin{equation}
\rho_\chi(r) \propto
\frac{e^{-2m_\chi r}}{r^2}.
\label{eq:rho-halo}
\end{equation}

The enclosed halo mass within radius $r$ is
\begin{equation}
M_\chi(<r) = 4\pi\int_0^r \rho_\chi(r')\,r'^2\,dr'
\propto \int_0^r e^{-2m_\chi r'}\,dr'
= \frac{1 - e^{-2m_\chi r}}{2m_\chi}.
\label{eq:M-enclosed}
\end{equation}

The halo contribution to the circular velocity follows from
$V_{\mathrm{halo}}^2 = GM_\chi(<r)/r$:
\begin{equation}
V_{\mathrm{halo}}^2(r) = B\,X_h(r),
\qquad
X_h(r) \equiv \frac{1 - e^{-2m r}}{r},
\label{eq:Vhalo}
\end{equation}
where $m \equiv m_\chi$ and $B$ is a phenomenological amplitude
absorbing the coupling constants and total halo mass scale.
The microscopic construction motivates the halo kernel $X_h$ and
its characteristic length scale $r_c = 1/m$, while the overall
amplitude $B$ is treated as a free parameter at galactic scales.
Future work deriving $B$ from the gap equation and baryon--condensate
coupling would provide a fully predictive model.

The kernel $X_h(r)$ has two limits:
\begin{itemize}
\item $r \ll r_c \equiv 1/m$:\;
  $X_h \approx 2m$, so that
  $V_{\mathrm{halo}}^2 \to 2mB = \mathrm{const}$
  (an approximately constant velocity plateau in the inner region),
\item $r \gg r_c$:\;
  $X_h \approx 1/r$ (Keplerian decline).
\end{itemize}

Note that while $V_{\mathrm{halo}}^2$ approaches a constant at
small $r$, the underlying density profile~\eqref{eq:rho-halo}
scales as $r^{-2}$ and does not possess a finite central density.
The inner velocity plateau is therefore a kinematic feature of the
enclosed mass profile rather than a density core in the usual sense.

\emph{Status of the mass scale.}  The identification of the
scalar-channel composite with a kpc-range field requires
$m_\chi \sim 10^{-26}$~eV, more than fifty orders of magnitude below the
substrate cutoff.  A companion lattice measurement of the
scalar-channel pole in the same framework finds instead
$m_\sigma \simeq 0.83 \times 2m_{\rm dyn}$, i.e.\ a composite at the
cutoff scale, with the ratio approaching the continuum
Nambu--Jona-Lasinio value $m_\sigma = 2m_{\rm dyn}$ in the chiral
limit.  The construction of this section is therefore retained as
motivation for the \emph{form} of the kernel $X_h$, not as a derivation
of its scale; the microscopic status of the mediator, and the
framework's actual light degree of freedom, are taken up in
Section~\ref{sec:a0}.


% ===================================================================
\section{Fitting Methodology}
\label{sec:method}
% ===================================================================

\subsection{Data}

We use the mass model data of \cite{Lelli:2016zqa}, providing
$V_{\mathrm{obs}}$, $\delta V$, $V_{\mathrm{gas}}$,
$V_{\mathrm{disk}}$, and $V_{\mathrm{bul}}$ for 175 galaxies.

\subsection{Rotation curve model}

\begin{equation}
V_{\mathrm{model}}^2(r) = V_{\mathrm{gas}}^2
+ \mu\,(V_{\mathrm{disk}}^2 + V_{\mathrm{bul}}^2)
+ B\,X_h(r),
\label{eq:Vmodel}
\end{equation}
where $\mu = \Upsilon_*/\Upsilon_{*,\mathrm{fid}}$ rescales the
stellar mass-to-light ratio.

\subsection{Two-stage fitting procedure}

\textbf{Stage 1: Parameter estimation in $V^2$ space.}
For fixed $m$, the parameters $(\mu, B)$ are obtained by weighted
least squares applied to
$V_{\mathrm{obs}}^2 - V_{\mathrm{gas}}^2 = \mu\,S_*(r) + B\,X_h(r)$,
with weights $w_i = 1/\sigma_{V,i}^2$ and
$\sigma_{V,i} = \max(\delta V_i, \sigma_f)$.
This approach exploits the $V^2$-linearity of the model and is
standard in the SPARC literature
\cite{Lelli:2016zqa,Li:2020iib,Lelli:2017vgz}.

\textbf{Stage 2: Model evaluation in $V$ space.}
Fit quality is assessed via
$\chi^2 = \sum_i [(V_{\mathrm{obs},i} - V_{\mathrm{model},i})/\sigma_{V,i}]^2$
and the full Gaussian log-likelihood including the $\log\sigma$ term.
Model comparison uses AIC with dynamic effective parameter count.

\textbf{Parameter ranges.}
The stellar rescaling is restricted to $\mu \geq 0$ (unbounded
above) and the halo amplitude to $B \geq 0$; the weighted
least-squares solution is projected onto these bounds, and
solutions at a boundary are reported as such.
In practice only two galaxies fit with $\mu > 10$ (F565-V2 at
$\mu = 13$ and UGCA444 at $\mu = 77$, the latter a poor fit);
all aggregate statistics are medians and are insensitive to them.
The unbounded range on $\mu$ is deliberately uninformative, so that
the analysis tests the halo kernel shape without imposing a
stellar-population prior; the consequences of this choice for
the $\mu$ distribution are discussed in
Section~\ref{sec:results}, and its decisive consequences for the
finite-$r_c$ population in Section~\ref{sec:ml-prior}.

\textbf{Configuration.}
The scalar mass $m$ is scanned over $[0.00125, 0.5]$~kpc$^{-1}$
($r_c \in [2, 800]$~kpc), with 140 log-spaced points.
For each galaxy, nine configurations of inner-radius mask
($r_{\min} \in \{0, 2, 3\}$~kpc) and error floor
($\sigma_f \in \{3, 4, 6\}$~km/s) are explored; the best
configuration is selected by AIC.
This systematic exploration accommodates the varying data quality
and radial coverage across the heterogeneous SPARC sample.
Since the number of configurations is small (nine) and fixed
across all galaxies, the effective trial factor is modest;
the robustness analysis in Section~\ref{sec:robust} confirms
that core findings are stable under fixed-configuration fitting.

\textbf{Boundary handling.}
When $m$ hits the lower grid edge, a 90\% CL lower limit on $r_c$
is reported via profile likelihood
($\Delta(-2\ln\mathcal{L}) = 2.71$).


% ===================================================================
\section{Results}
\label{sec:results}
% ===================================================================

\subsection{Overview}

\begin{table}[h]
\centering
\begin{tabular}{lc}
\toprule
Quantity & Value \\
\midrule
Total galaxies fitted & 175 \\
Finite $r_c$ & 101 (58\%) \\
Boundary $r_c$ & 74 (42\%) \\
Median $r_c$ (finite) & 15.8~kpc \\
10--90th percentile $r_c$ & 2.0--185~kpc$^{\dagger}$ \\
Median $\chi^2/\mathrm{pt}$ & 1.08 \\
\multicolumn{2}{l}{\footnotesize $^{\dagger}$lower edge grid-censored;
extended-grid 10th pct $= 0.16$~kpc (Sec.~\ref{sec:robust})} \\
$\chi^2/\mathrm{pt} < 1$ & 81 (46\%) \\
\midrule
Median $\mu$ & 0.095 \\
IQR of $\mu$ & 0.00--1.25 \\
$\mu = 0$ (boundary) & 81 (46\%) \\
\bottomrule
\end{tabular}
\caption{Summary of Yukawa halo fits to 175 SPARC galaxies.
The high fraction of $\mu = 0$ solutions indicates that many
galaxies are well-fitted with negligible stellar contribution
to the rotation curve, a feature also observed in other halo
model analyses of the SPARC sample \cite{Li:2020iib}.}
\label{tab:overview}
\end{table}

The prevalence of $\mu = 0$ solutions deserves comment.
The inner velocity plateau of the Yukawa kernel
($V_{\mathrm{halo}}^2 \to 2mB$ as $r \to 0$;
Section~\ref{sec:micro}) is partially degenerate with the stellar
contribution over the limited radial range of many rotation
curves, so that an uninformative prior on $\mu$ allows the halo
component to absorb part of the stellar signal.
Similar $\Upsilon_*$ boundary behavior is observed in
unconstrained halo fits to SPARC
\cite{Li:2020iib}, and is conventionally addressed by imposing
a lognormal population prior on $\Upsilon_*$.
The fits reported here should therefore be read as a test of the
halo kernel \emph{shape} under maximally agnostic stellar
masses; an analysis with an informative $\Upsilon_*$ prior is an
important consistency check and is addressed in
Section~\ref{sec:robust}.

\subsection{Cutoff radius distribution}

The results of this and the following subsections characterize the
free-$\mu$ analysis; their prior sensitivity is quantified in
Section~\ref{sec:ml-prior}.
Among the 101 galaxies with finite $r_c$, the distribution
is broad but shows structure:

\begin{center}
\begin{tabular}{cc}
\toprule
$r_c$ range (kpc) & Count \\
\midrule
$[0, 5)$ & 26 \\
$[5, 10)$ & 15 \\
$[10, 20)$ & 11 \\
$[20, 30)$ & 8 \\
$[30, 50)$ & 7 \\
$[50, 100)$ & 13 \\
$[100, 200)$ & 11 \\
$[200, 800)$ & 10 \\
\bottomrule
\end{tabular}
\end{center}

The emergent scalar mass $m_\chi \sim 0.06$~kpc$^{-1}$ at the
median lies in the ultralight range ($m_\chi \sim 10^{-27}$~eV).
This is significantly below the canonical fuzzy dark matter mass
scale $m_{\mathrm{FDM}} \sim 10^{-22}$~eV
\cite{Hu:2000ke,Hui:2016ltb}, indicating that the emergent scalar
field occupies a distinct region of ultralight parameter space.
We emphasize that the standard fuzzy-dark-matter constraints from
small-scale structure---most notably the Lyman-$\alpha$ forest,
which bounds the mass of a \emph{relic} wave-dark-matter
component---do not apply directly here.  In the present framework
the halo is a static field \emph{sourced by baryons}
(Eq.~\ref{eq:scalar-field}) rather than a virialized relic of an
independent condensate, so $m_\chi$ sets a per-galaxy length scale
rather than a free-streaming or de~Broglie scale.  As we argue in
Section~\ref{sec:a0}, the variation of this length across galaxies
is itself best read as reflecting an underlying acceleration scale
rather than a fixed field mass---though, as the direct test of
Section~\ref{sec:ac-test} shows, $r_c$ marks where the fixed-range
kernel saturates against the baryon distribution, not where the
acceleration crosses $a_0$.

\subsection{Test of a universal scalar mass}
\label{sec:universal-m}

If the scalar mass $m_\chi$ is a fundamental microscopic constant,
all galaxies should be well-described by a single value of $m$.
We test this by fitting all 175 galaxies with $m$ fixed at a
common value, optimizing only $(\mu, B)$ per galaxy.

The best universal value is $m \approx 0.005$~kpc$^{-1}$
($r_c \approx 200$~kpc), which yields a median
$\chi^2/\mathrm{pt} = 1.31$---a 22\% degradation relative to
the per-galaxy fit (1.08).
At this universal $m$, 62\% of galaxies achieve
$\chi^2/\mathrm{pt} < 2$.
In contrast, fixing $m$ at the per-galaxy median value of
$0.063$~kpc$^{-1}$ ($r_c = 15.8$~kpc) produces a median
$\chi^2/\mathrm{pt} = 4.35$, indicating that a ``typical'' $m$
does not adequately describe the majority of the sample.

The per-galaxy best-fit $m$ values span approximately 1.4~decades
(standard deviation 0.72~dex in $\log_{10}m$), confirming that
the data do not support a single universal scalar mass within the
present linearized, point-source framework.

This finding admits three interpretations that are not mutually
exclusive:
\begin{enumerate}
\item The effective scalar mass may be environment-dependent,
  varying with local halo density, baryonic concentration,
  or formation history. Such dependence could arise from the
  nonlinear $\chi^3$ term in Eq.~\eqref{eq:scalar-field},
  which shifts the effective mass in regions of high condensate
  amplitude.
\item The linearized point-source approximation
  (Eq.~\ref{eq:yukawa-point}) may be insufficient: the full
  solution of the nonlinear field equation with spatially
  extended baryonic sources could produce an effectively
  variable cutoff scale even with a single microscopic $m_\chi$.
\item Most pointedly, the spread in $m$ is not random but trends
  with the galaxy's characteristic velocity
  (Section~\ref{sec:subsample}): the data select a length that
  \emph{scales with the system}, as expected if the governing
  physics is a universal \emph{acceleration} scale rather than a
  universal length.  A direct test of the strongest form of this
  reading---that the per-galaxy $r_c$ marks the radius at which the
  local acceleration crosses a fixed threshold $a_0$---is performed
  in Section~\ref{sec:ac-test} and \emph{fails}: the characteristic
  accelerations at $r_c$ are not universal.  The acceleration-scale
  interpretation survives only at the pointwise level, with the
  Yukawa $r_c$ an artifact of forcing a fixed-range kernel onto an
  acceleration-dependent response.  We develop this distinction,
  and its origin in the framework's dark-energy sector, in
  Sections~\ref{sec:a0} and~\ref{sec:ac-test}.
\end{enumerate}

Distinguishing these scenarios requires either a nonlinear
numerical solution of Eq.~\eqref{eq:scalar-field} with realistic
source profiles, or a hierarchical Bayesian analysis testing
a narrow distribution of $m$ values.
Both are left for future work.
We note that this test, like all of Section~\ref{sec:results}, is
performed within the free-$\mu$ analysis; under the physical
stellar prior of Section~\ref{sec:ml-prior} the finite-$r_c$
population itself largely vanishes, which supersedes rather than
rescues the universal-mass hypothesis---no single $m$ describes
the sample in either treatment.

% --- Figures ---
\begin{figure*}[ht]
\centering
\includegraphics[width=\textwidth]{fig1_rotation_curves.pdf}
\caption{Representative Yukawa halo fits to four SPARC galaxies.
\textit{Top left}: NGC2998, a massive spiral with finite
$r_c = 219.5$~kpc and Yukawa preferred over NFW
($\Delta\mathrm{AIC} = -2.2$).
\textit{Top right}: NGC4138, with well-constrained
$r_c = 11.2$~kpc.
\textit{Bottom left}: DDO154, a dwarf galaxy classified as a
boundary case ($r_c \geq 800$~kpc).
\textit{Bottom right}: NGC6503, a boundary case with excellent
fit quality ($\chi^2/\mathrm{pt} = 0.18$).
In each panel, black points with error bars show observed
velocities, the red curve is the total model, the green dashed
curve is the Yukawa halo contribution, and the shaded regions
show gas (cyan) and stellar (blue) components.}
\label{fig:fits}
\end{figure*}

\begin{figure}[ht]
\centering
\includegraphics[width=0.9\columnwidth]{fig2_rc_distribution.pdf}
\caption{Distribution of finite cutoff radii $r_c = 1/m$ for
the 101 galaxies with well-constrained fits.
The red dashed line marks the median $r_c = 15.8$~kpc.
The distribution is broad; the apparent pile-up at the 2~kpc grid
edge is partly censoring---on an extended grid 16 of the 26
galaxies in $[2,5)$~kpc migrate below 2~kpc
(Section~\ref{sec:robust}), widening the span to
$[0.13, 734]$~kpc.}
\label{fig:rc-dist}
\end{figure}

\begin{figure}[ht]
\centering
\includegraphics[width=0.9\columnwidth]{fig3_dAIC_distribution.pdf}
\caption{Distribution of $\Delta\mathrm{AIC} =
\mathrm{AIC}_{\mathrm{Yukawa}} - \mathrm{AIC}_{\mathrm{NFW}}$
for all 175 galaxies.
Negative values indicate Yukawa preference; positive values
indicate NFW preference.
The red dashed line marks the median $\Delta\mathrm{AIC} = +6.6$.
The green dotted lines mark $\pm 2$, the conventional threshold
for meaningful model distinction.}
\label{fig:daic}
\end{figure}


% ===================================================================
\section{Comparison with the NFW Profile}
\label{sec:nfw}
% ===================================================================

\subsection{NFW fitting}

The NFW profile is fitted to all 175 galaxies using identical
methodology, with $r_s$ scanned over $[0.5, 200]$~kpc (80 points).

\subsection{Overall comparison}

\begin{table}[h]
\centering
\begin{tabular}{lcc}
\toprule
Category & Count & Fraction \\
\midrule
Yukawa preferred ($\Delta\mathrm{AIC} < -2$) & 10 & 6\% \\
Comparable ($|\Delta\mathrm{AIC}| \leq 2$) & 48 & 27\% \\
NFW preferred ($\Delta\mathrm{AIC} > 2$) & 117 & 67\% \\
\midrule
Median $\Delta\mathrm{AIC}$ & \multicolumn{2}{c}{$+6.6$} \\
\bottomrule
\end{tabular}
\caption{AIC comparison: Yukawa vs NFW.
$\Delta\mathrm{AIC} = \mathrm{AIC}_{\mathrm{Yuk}} - \mathrm{AIC}_{\mathrm{NFW}}$.}
\label{tab:nfw}
\end{table}

While NFW is generally preferred, the Yukawa profile achieves
comparable or superior performance in 33\% of galaxies (58 systems).
The median $\Delta\mathrm{AIC} = +6.6$ represents a moderate
preference, below the threshold for strong evidence ($\Delta\mathrm{AIC} > 10$).
$\Delta\mathrm{AIC}$ is evaluated with both kernels fitted on the
same data cut and error floor---the $(r_{\min}, \sigma_f)$
configuration selected by the Yukawa fit---since AIC values are
comparable only on identical data with an identical noise model.
Selecting each kernel's own AIC-optimal configuration instead
shifts the median to $+9.5$, further in NFW's favour, moving four
galaxies across the $|\Delta\mathrm{AIC}| = 2$ boundary without
altering any conclusion.

\subsection{Subsample analysis by galaxy mass}
\label{sec:subsample}

To investigate whether the Yukawa profile performs differently
across galaxy types, we partition the sample by
maximum observed velocity $V_{\max}$:

\begin{table}[h]
\centering
\begin{tabular}{lccc}
\toprule
Subsample & $N$ & Median $\Delta$AIC & Yuk.\ comp.\ or better \\
\midrule
$V_{\max} < 80$ km/s (dwarf/LSB) & 59 & $+2.8$ & 42\% \\
$80 \leq V_{\max} < 150$ km/s & 60 & $+13.0$ & 22\% \\
$V_{\max} \geq 150$ km/s (massive) & 56 & $+11.3$ & 36\% \\
\bottomrule
\end{tabular}
\caption{NFW comparison by galaxy mass subsample.
``Comp.\ or better'' means $\Delta\mathrm{AIC} \leq 2$.}
\label{tab:subsample}
\end{table}

Two features are notable.
First, the dwarf/LSB subsample shows the smallest median
$\Delta\mathrm{AIC} = +2.8$, suggesting a reduced NFW advantage
in low-mass systems.
The inner velocity plateau of the Yukawa profile
(Section~\ref{sec:micro}) may contribute to this improved
relative performance, though a detailed investigation of
inner density slopes is beyond the scope of this paper.
Second, 7 of the 10 galaxies where Yukawa is decisively preferred
($\Delta\mathrm{AIC} < -2$) are massive spirals
($V_{\max} \geq 150$~km/s), suggesting that the Yukawa cutoff
scale $r_c$ may be particularly well-suited to the halo structure
of these systems.

\subsection{Context}

Several considerations inform the interpretation of this comparison:

\begin{enumerate}
\item The NFW profile is an empirical fit from $N$-body simulations,
  refined over three decades.
  The Yukawa profile is derived from a microscopic construction
  with no optimization for rotation curve phenomenology.
  Achieving performance within $\Delta\mathrm{AIC} \sim 7$ of NFW
  is a nontrivial result for a microscopically motivated profile.

\item The comparison is at equal per-galaxy parameter count: the
  NFW inner and outer slopes are fixed ($-1$ and $-3$), so both
  kernels carry exactly two per-galaxy parameters (a scale and an
  amplitude).  The AIC difference therefore reflects fit quality
  alone, not functional flexibility.

\item The Yukawa profile produces an inner velocity plateau
  ($V^2 \to \mathrm{const}$ as $r \to 0$), providing a qualitatively
  different inner structure from the NFW profile.
  The reduced NFW advantage in dwarf galaxies is consistent with
  this structural difference, though a full density-profile
  comparison is deferred to future work.
\end{enumerate}


% ===================================================================
\section{Robustness Analysis}
\label{sec:robust}
% ===================================================================

\subsection{Fitting space: $V^2$ vs.\ $V$}

Performing both parameter estimation and scoring entirely in
$V$ space (via iterative nonlinear optimization) yields
qualitatively consistent results: the finite/boundary split
shifts from 101/74 to 95/80, and $\chi^2/\mathrm{pt}$ increases
from 1.08 to 1.32.
The NFW comparison remains similar.
The $V^2$ approach is adopted as primary because it is standard
in the SPARC literature and exploits the physical linearity of
the circular velocity equation.

\subsection{Error treatment}

Replacing $\sigma_V = \max(\delta V, \sigma_f)$ with quadrature
addition $\sigma_V = \sqrt{\delta V^2 + \sigma_f^2}$ produces
negligible changes in the aggregate statistics.

\subsection{Grid boundaries}

Varying $m_{\min}$ from 0.001 to 0.002~kpc$^{-1}$
($r_{c,\max} = 500$--$1000$~kpc) changes the boundary fraction
by $\lesssim 5\%$, confirming the 58\%/42\% split is not a
grid artifact.

\subsection{Stellar mass-to-light prior}
\label{sec:ml-prior}

As noted in Section~\ref{sec:results}, the uninformative range on
$\mu$ permits solutions with vanishing stellar contribution.
To test whether the core findings depend on this freedom, we
repeat the analysis imposing a lognormal prior on the stellar
mass-to-light rescaling, centered on $\mu = 1$
(i.e.\ the fiducial $\Upsilon_*$) with a scatter of 0.1~dex,
following standard practice in SPARC halo fitting
\cite{Li:2020iib}, with the same prior applied to both kernels.
Two implementations bracket the result.
In the first, $\mu$ is selected by the penalized (maximum a
posteriori) Gaussian likelihood in velocity space over a
0.02-dex grid, with $B$ solved at each $\mu$ by the same
$V^2$-space weighted least squares as the fiducial pipeline.
In the second, $\mu$ is fixed at unity and the machinery is
otherwise \emph{identical} to the fiducial pipeline, removing any
prior-strength ambiguity.

The result is a qualitative change, not a perturbation.
Under the lognormal prior, the finite/boundary split becomes
12/163 (7\%/93\%), the median finite $r_c$ becomes 6.4~kpc, the
median $\chi^2/\mathrm{pt}$ becomes 1.14, and the median
$\Delta\mathrm{AIC}$ relative to NFW (fitted under the same
prior) becomes $+7.7$.
With $\mu$ fixed at unity, the split becomes 4/171 (2\%/98\%),
the median $\chi^2/\mathrm{pt}$ rises to 1.48, the median
$\Delta\mathrm{AIC}$ is $+6.6$, and the number of
Yukawa-preferred galaxies falls from 10 to zero.
The galaxies whose finite $r_c$ survives the physical
normalization (DDO168, NGC2366, UGC05764) are gas-dominated
dwarfs, for which the stellar normalization is immaterial by
construction.

The qualitative conclusions of the free-$\mu$ analysis are
therefore \emph{not} preserved: the finite-$r_c$ population is a
creature of the unconstrained stellar freedom---the kernel
achieves its fiducial success substantially by suppressing the
stellar contribution and absorbing it into the halo shape---and
under a physical $\Upsilon_*$ the sample almost uniformly prefers
the unscreened ($r_c \to \infty$) limit.
We record one further sensitivity uncovered by this test: with
the velocity-space MAP estimator and \emph{no} prior, the split
is 44/131 (25\%/75\%), so the finite/boundary classification also
depends appreciably on the estimator itself.
Both facts point the same way: $r_c$ is not a robustly
constrained physical scale in this sample, which is precisely
what is expected if the Yukawa kernel is the wrong functional
form and its per-galaxy freedom is absorbing missing
acceleration-dependent physics
(Sections~\ref{sec:a0} and~\ref{sec:ac-test}).
The model-comparison verdict---NFW moderately preferred---is
unchanged in every variant.



\subsection{Lower-boundary censoring of $r_c$}

The fiducial grid spans $r_c \in [2, 800]$~kpc, and the upper
boundary is treated as censored (the 74 boundary systems).  The
lower boundary deserves---and on inspection requires---the same
treatment: the $[2,5)$~kpc bin is the most populated (26 galaxies)
and the finite-$r_c$ 10th percentile sits exactly on the 2~kpc grid
edge.  Extending the grid to $r_c \in [0.25, 800]$~kpc (the
fiducial grid an exact subset; all 175 galaxies rerun) shows the
pile-up to be partly a wall artifact: 16 of the 26 migrate below
2~kpc (15 of them had sat exactly on the edge), with median
$\Delta\mathrm{AIC} = -1.4$ and up to $-28.5$ (UGC04278); 8 of
the 16 run to the new grid edge at $0.13$~kpc, preferring an
ever-smaller cutoff---no preferred scale at all, the Yukawa kernel
degenerating toward a point-mass term.  The finite-$r_c$ 10th
percentile drops from $2.0$ to $0.16$~kpc and the span widens to
$[0.13, 734]$~kpc.  The correction \emph{strengthens} the
non-universality failure mode: the true spread of the preferred
range is wider than the fiducial grid could express.  Statements
involving small-$r_c$ systems elsewhere in this paper (notably the
$r_c$--$V_{\max}$ scalings of Section~\ref{sec:ac-test}) carry
this censoring caveat.  (Scripts and per-galaxy table:
\texttt{rc\_floor.py}, \texttt{results/rc\_floor.csv}.)

\subsection{Configuration selection}

The primary analysis selects the best of nine
$(r_{\min}, \sigma_f)$ configurations per galaxy.
To verify that this selection does not introduce significant
bias, we also perform a fixed-configuration analysis using
$r_{\min} = 0$~kpc and $\sigma_f = 6$~km/s uniformly for all
galaxies.
Under fixed configuration, the finite/boundary split shifts to
97/78 (55\%/45\%), the median $r_c$ shifts to 7.3~kpc, and the
median $\chi^2/\mathrm{pt}$ increases to 1.54.
The core qualitative findings---a substantial fraction of
finite-$r_c$ galaxies, a broad $r_c$ distribution, and
moderate NFW preference---are preserved.
The multi-configuration analysis is adopted as primary because
it better accommodates the heterogeneous data quality and radial
coverage of the SPARC sample.


% ===================================================================
\section{Discussion}
\label{sec:discuss}
% ===================================================================

\subsection{Relation to previous Yukawa studies}

Previous applications of Yukawa potentials to galaxy rotation
curves \cite{deAlmeida:2018kwq,Henrichs:2021bnd,Hassan:2025yuk}
introduce the Yukawa correction as a phenomenological fifth force
arising from $f(R)$ or scalar-tensor modifications of gravity.
In those frameworks, the Yukawa coupling $\beta$ and range $\lambda$
are free parameters with no specified microphysical origin.

In contrast, the present work derives the halo functional form
from a definite microscopic construction (NJL-type four-fermion
interaction on a lattice), where the halo kernel $X_h(r)$ and its
characteristic length scale $r_c = 1/m$ are consequences of
Hubbard--Stratonovich bosonization rather than a phenomenological
ansatz.
The amplitude $B$ and mass-to-light rescaling $\mu$ remain
phenomenological parameters at the galactic level; the microscopic
construction motivates the kernel shape and its length scale, but
does not yet predict the overall normalization.

It is important to note that $B$ is not an arbitrary dark halo mass
analogous to the NFW $M_{200}$.
In the microscopic theory, the condensate field $\chi$ is directly
sourced by the baryonic density $J_B$
(Eq.~\ref{eq:scalar-field}), so the halo amplitude is in
principle determined by the baryon--condensate coupling $\mu_c$
and the total baryonic mass $M_B$.
This baryon-sourced structure means that the halo normalization
is not independent of the visible matter distribution---a feature
that may eventually provide a physical basis for the empirical
correlations between baryonic and dynamical masses observed in
galaxies \cite{McGaugh:2016leg,Lelli:2017vgz}, though the present
paper does not yet derive this connection quantitatively.

Conceptually, a baryon-sourced halo places the model between
particle dark matter and modified gravity: the dynamical mass
tracks the visible matter through Eq.~\eqref{eq:scalar-field}, as in
modified-gravity scenarios, yet it arises from a definite
microscopic field rather than a postulated force law.  This
positioning carries the corresponding obligations.  The model must
ultimately confront the tests that distinguish a sourced force from
collisionless matter---most sharply systems in which the
gravitational and baryonic centroids are spatially displaced (e.g.\
merging clusters such as the Bullet Cluster), together with
cluster-scale dynamics and the acoustic peaks of the cosmic
microwave background.  These probes lie beyond the galactic
rotation curves examined here and are deferred to future work.

In the present analysis, $B$ is nevertheless treated as a free
parameter.
This is a deliberate choice: by decoupling the test of the
\emph{spatial geometry} of the halo kernel $X_h(r)$ from the
\emph{absolute calibration} of the microscopic coupling $\mu_c$,
we avoid conflating uncertainties in stellar mass-to-light ratios
with uncertainties in the theoretical normalization.
The present paper therefore tests a specific and falsifiable
prediction---the Yukawa radial shape and its characteristic
cutoff scale $r_c$---while deferring the absolute normalization
to future work in which $B$ would be derived from the gap equation
and baryon--condensate coupling.

\subsection{An emergent acceleration scale}
\label{sec:a0}

The non-universality of the fitted mass (Section~\ref{sec:universal-m})
and its trend with galaxy velocity (Section~\ref{sec:subsample}) point
beyond the fixed-range picture. Two considerations sharpen this.

\emph{A constraint.} Read literally as a fixed-range field of Compton
wavelength $r_c \sim 15$~kpc that is sourced by baryons and acts on
baryons, the halo is a Yukawa fifth force. For a matter-coupled
scalar the direct exchange force dominates the self-gravity of the
sourced cloud by a factor $\sim (G M_B/R_B c^2)^{-1} \gg 1$, so the
rotation-curve effect is carried by the force, not by an extra
gravitating mass. To reproduce the missing acceleration this force
must be of gravitational strength on galactic scales; since its range
vastly exceeds the Solar System, the same force would be essentially
unscreened there, producing a post-Newtonian deviation
$|\gamma-1|\sim\mathcal{O}(1)$ in conflict with the Cassini bound
$|\gamma-1| < 2.3\times10^{-5}$. A fixed-range, universal Yukawa
interpretation is therefore disfavoured---independently of the
fitting results, and consistent with the data's own preference
against a universal mass.

\emph{A resolution.} The empirical pattern---a length that scales with
the system, equivalently a fixed acceleration scale---is the hallmark
of the radial acceleration relation
\cite{McGaugh:2016leg,Lelli:2017vgz}. Both the constraint and the
pattern are met if the scalar response is \emph{acceleration
dependent} rather than fixed-range: a non-canonical (k-mouflage)
kinetic sector, $\mathcal{L} = \Lambda_{\rm DE}^4\,
K\!\big((\partial\chi)^2/\Lambda_{\rm DE}^4\big) - \chi\,
T^\mu_{\ \mu}/M_{\rm Pl}$, screens the force where the local
acceleration is large (the Solar System, $a \gg a_0$) and restores it
where the acceleration is small (galaxy outskirts, $a \lesssim a_0$).
Absence of ghost and gradient instabilities ($K' > 0$, $c_s^2 \le 1$)
selects the enhancing, MOND-like branch; the opposite sign is not an
alternative phenomenology but a pathology.

The acceleration scale is not free. In this framework the infrared is
dominated by the dark-energy density of the same condensate,
$\rho_{\rm DE}\sim\Lambda_{\rm DE}^4$, and the crossover acceleration
where the kinetic sector turns over,
$(\partial\chi)^2\sim\Lambda_{\rm DE}^4$, is
\begin{equation}
a_0 \;\sim\; \frac{\Lambda_{\rm DE}^2}{M_{\rm Pl}^{\rm (full)}}
\;\approx\; 2\times10^{-10}\ \mathrm{m\,s^{-2}}
\;\sim\; \frac{cH_0}{2\pi},
\label{eq:a0}
\end{equation}
within a factor of two of the rotation-curve value
$a_0^{\rm obs}\approx1.2\times10^{-10}\ \mathrm{m\,s^{-2}}$.
Equation~\eqref{eq:a0} ties the galactic acceleration scale to the
cosmological one, giving a structural reading of the long-noted
coincidence $a_0\sim cH_0\sim c^2\sqrt{\Lambda}$: the dark-matter and
dark-energy phenomenologies share a single scale because they are
different infrared aspects of one condensate.

\emph{The mediator's identity.}  The microphysics independently
excludes the fixed-mass scalar as the galactic mediator: the
companion lattice measurement places the scalar-channel composite
at the dynamical fermion scale ($m_\sigma \simeq 0.83 \times
2m_{\rm dyn}$---a microscopic mass, $m_\sigma \gg \mathrm{kpc}^{-1}$,
at the Planck-scale operating point---
approaching the continuum value $2m_{\rm dyn}$ in the chiral
limit; Section~\ref{sec:micro}), so the $\sigma$ integrates out
more than fifty orders of magnitude above any galactic scale.
The framework's light degree of freedom is instead the
condensate's \emph{angular pseudo-Goldstone}, whose potential is
protected to be exponentially flat.  Its coupling route, however,
is constrained by two results computed since the fiducial analysis
(companion dark-energy paper).  First, the superfluid-phonon
($P(X)$) reading is closed by a two-body theorem: the mean-field
equation of state of a contact pairing channel is
$P \propto (\mu-\mu_c)^2$ at onset for any coupling strength,
never the $(\mu-\mu_c)^{3/2}$ a MOND-like force law requires---the
structure superfluid dark matter models
\cite{Berezhiani:2015bqa} must \emph{postulate}.  Second, parity
forbids the pseudo-Goldstone any linear coupling to unpolarized
matter, and the disformal coupling it does acquire at loop level
carries Planck-scale coefficients, suppressed at galactic gradients
by $(b/m_{\rm dyn})^2 \sim 10^{-122}$.
The remaining candidate was the gradient response of the
topological (monopole) sector that was, at the time, also the
candidate generator of the dark-energy scale---a candidacy since
terminated by direct measurement in the companion paper (its
susceptibility sits at the cutoff scale, not the meV scale).  That
computation, pre-registered in the companion paper, has now been
executed at leading order: the dilute gas acts on the angular mode
through the periodic potential it generates, so the response is the
sine-Gordon soliton lattice, solvable exactly.  The verdict is sharp
on both sides.  The crossover lands at precisely the required scale,
$X \sim \Lambda_{\rm DE}^4$: the scale selection works.  But the
low-$X$ response is $L \propto X^{1/2}$ (wall-lattice pinning),
never the $X^{3/2}$ of a MOND force law, and, independently, nothing
sources galactic $\theta$-gradients in the first place
($\theta = \mathrm{const}$ solves the field equation in any static
matter background under the allowed disformal coupling).  In the
present dilute-monopole realization, the same topology that
generates the meV crossover does so through a local periodic
potential, and a potential yields walls, not MOND.  The
emergent-MOND reading of this framework therefore fails at leading
order, and Eq.~\eqref{eq:a0} stands as an unexplained but robust
structural coincidence.

We are explicit about status. The scale, Eq.~\eqref{eq:a0}, and the
direction (MOND-like, by consistency) are robust; a first-principles
demonstration that the condensate generates the requisite k-mouflage
kinetic function---with an $\mathcal{O}(1)$ coefficient of the correct
sign---remains to be completed and is the principal theoretical task
left open here. A companion analysis develops the dark-energy sector
that fixes $\Lambda_{\rm DE}$.

\subsection{A direct test: characteristic accelerations at the
fitted cutoff radii}
\label{sec:ac-test}

The strongest form of the acceleration-scale reading makes a
falsifiable prediction: if the per-galaxy $r_c$ marks the radius at
which the local acceleration crosses a universal threshold, then
$a_c \equiv V_{\rm obs}^2(r_c)/r_c$ should cluster at a single value
$\sim a_0$, independently of galaxy properties.  We test this with
criteria fixed in advance: clustering requires a scatter
$\lesssim 0.3$~dex in $\log_{10} a_c$ and no significant residual
trend with $V_{\max}$.  To avoid extrapolation, only galaxies whose
$r_c$ lies within the radial coverage of the data are evaluated; in
the fiducial run this retains 33 of the 101 finite-$r_c$ galaxies,
a selection (toward smaller $r_c$) that we note explicitly.

The test fails (Fig.~\ref{fig:ac-test}).  The retained accelerations
have median $a_c = 2.7\times10^{-11}\ \mathrm{m\,s^{-2}} \approx
0.22\,a_0^{\rm obs}$ with a scatter of $0.34$~dex, and---decisively---a
strong systematic trend with velocity:
$d\log a_c/d\log V_{\max} = 1.09 \pm 0.17$
(Spearman $\rho = 0.72$, $p = 3\times10^{-6}$).
The verdict is robust across methodology: an alternative run using
the extended-source convolution kernel and $V^2$-space error
propagation, with 100 galaxies inside coverage, gives a slope of
$1.33 \pm 0.10$ and a scatter of $0.44$~dex.
A constant characteristic acceleration would require
$r_c \propto V_{\max}^{2}$; the measured $r_c$--$V_{\max}$ scaling
is not only shallower but is itself unstable across samples:
$0.51 \pm 0.24$ on the 33-galaxy subsample above,
$0.82 \pm 0.15$ on the 156 finite-$r_c$ galaxies of the
extended-source run, and $1.39 \pm 0.22$ on the 101 finite-$r_c$
galaxies of the fiducial run.  No single exponent describes it, and
the instability is itself diagnostic: $r_c$ is not a well-defined
physical scale of the sample---a further symptom, alongside the
non-universality of $m$, of a parametrization organizing itself
around the baryon distribution rather than around any physical
length.  The statement that survives every sample and methodology
is the one that matters: $a_c$ \emph{tracks} $V_{\max}$ (slopes
$1.09 \pm 0.17$ and $1.33 \pm 0.10$), far from the zero slope a
constant $a_c$ would require.

Two conclusions follow.  First, the per-galaxy $r_c$ of the Yukawa
parametrization is \emph{not} an acceleration-scale proxy: the
fixed-range kernel, forced onto an acceleration-dependent response,
selects a length governed by the baryon distribution rather than by
the $a_0$ crossing.  This is the second, independent symptom of the
wrong-kernel diagnosis (the first being the non-universality of $m$
itself, Section~\ref{sec:universal-m}).
Second, the failure does \emph{not} exclude an underlying universal
$a_0$ at the pointwise level, where it is in fact well established:
the radial acceleration relation in these same data, with stellar
mass-to-light ratios fixed at population-synthesis values, exhibits
$\sim 0.13$~dex scatter about a single
$a_0 \approx 1.2\times10^{-10}\ \mathrm{m\,s^{-2}}$
\cite{McGaugh:2016leg,Lelli:2017vgz}.
Our own fits cannot presently adjudicate this level: with the
deliberately uninformative unbounded $\mu \geq 0$ of
Section~\ref{sec:method}, the pointwise relation constructed from
the fitted $\mu$ values is substantially degraded (and in the
fiducial configuration destroyed), because unconstrained per-galaxy
mass-to-light rescalings absorb precisely the baryonic-acceleration
information the relation encodes.  The operational lesson for the
planned refit is therefore concrete: an acceleration-dependent
kernel with a \emph{single} fitted $a_0$, stellar mass-to-light
priors in place of free $\mu$, and evaluation in the pointwise
$g_{\rm obs}$--$g_{\rm bar}$ plane rather than through any derived
per-galaxy length.

\begin{figure}[ht]
\centering
\includegraphics[width=0.98\columnwidth]{fig4_ac_test.pdf}
\caption{Direct test of the $r_c$-as-transition-radius reading
(fiducial run, 33 galaxies with $r_c$ inside the radial coverage).
\textit{Left}: characteristic acceleration
$a_c = V_{\rm obs}^2(r_c)/r_c$ against $V_{\max}$; the fitted slope
$1.09\pm0.17$ (red) is inconsistent with the constant-$a_c$
prediction (gray dashed: $a_0^{\rm obs}$), which would require zero
slope.
\textit{Right}: distribution of $\log_{10} a_c$; scatter 0.34~dex.}
\label{fig:ac-test}
\end{figure}

This analysis is performed on the free-$\mu$ fiducial fits and
inherits their stellar freedom; under the physical prior of
Section~\ref{sec:ml-prior} the finite-$r_c$ population itself
largely vanishes, which supersedes rather than alters the
conclusion---in neither treatment does an $r_c$-level $a_0$
crossing exist.
Empirically, then, the refit of the SPARC sample with an
acceleration-dependent profile is the analysis this result demands;
it is performed in Section~\ref{sec:refit}.



\section{A single-acceleration-scale refit}
\label{sec:refit}

The three failure modes documented above prescribe the analysis
that must replace the fixed-range fit: an acceleration-dependent
kernel with a \emph{single} global scale, physical stellar
normalization, and no per-galaxy halo freedom.  This section
performs that refit and compares it, under strictly identical
methodology, with the Yukawa and NFW kernels.

\subsection{Kernel, nuisances, and likelihood}

For the acceleration-dependent response we adopt the empirical
radial-acceleration-relation (RAR) form
\cite{McGaugh:2016leg},
\begin{equation}
g_{\rm pred} \;=\; \frac{g_{\rm bar}}{1 - e^{-\sqrt{g_{\rm bar}/a_0}}},
\qquad
g_{\rm bar}(r) = \frac{V_{\rm gas}^2 + \mu\,(V_{\rm disk}^2 + V_{\rm bul}^2)}{r},
\label{eq:rar-kernel}
\end{equation}
with $a_0$ a single global parameter shared by all 175 galaxies.
Equation~\eqref{eq:rar-kernel} is a phenomenological stand-in for
the $K(X)$ response of Section~\ref{sec:a0}; its shape is an
empirical input, not a derivation, and replacing it with the
condensate-derived kinetic function remains the open theoretical
step.

Each galaxy carries three nuisance parameters with observational
priors: the stellar rescaling $\mu$ (lognormal, 0.1~dex about the
fiducial $\Upsilon_*$, as in Section~\ref{sec:ml-prior}), the
distance factor $d = D/D_0$ (Gaussian, width $e_D/D_0$ from the
SPARC catalog; $r \to d\,r$ and the flux-based
$V_{\rm bar}^2 \to d\,V_{\rm bar}^2$), and the inclination $i$
(Gaussian, width $e_i$; $V_{\rm obs} \to V_{\rm obs}\sin i_0/\sin i$)
\cite{Li:2018tdo}.
In place of the arbitrary error floor of the fiducial analysis, the
likelihood carries a single global fractional intrinsic scatter,
\begin{equation}
\sigma_{{\rm tot},i}^2 \;=\; (\sigma_{V,i}\,s)^2
 + (f_{\rm int}\,V_{{\rm obs},i}\,s)^2 ,
\end{equation}
where $s \equiv \sin i_0/\sin i$ is the inclination rescaling
applied to the observed velocities, so that the noise model
transforms with the data it describes.  The likelihood is fitted by
maximum likelihood (the Gaussian normalization terms
penalize inflating $f_{\rm int}$); the noise is attached to the
\emph{observed} velocities so that the error model is identical
across the models being compared.
The Yukawa and NFW kernels are fitted under exactly the same
nuisance structure and each with its own global $f_{\rm int}$,
retaining their per-galaxy scale and amplitude parameters.

\subsection{Results}

\begin{table}[t]
\centering
\begin{tabular}{lccccc}
\hline
Model & Global & Per-galaxy & $f_{\rm int}$ &
$\Delta\mathrm{AIC}$ & median $\chi^2/\mathrm{pt}$ \\
\hline
Single-$a_0$ (RAR) & $a_0$, $f_{\rm int}$ & --- & 5.0\% & $0$ & 0.83 \\
Yukawa & $f_{\rm int}$ & $m$, $B$ & 15.6\% & $+4859$ & 0.75 \\
NFW & $f_{\rm int}$ & $r_s$, $B$ & 3.4\% & $-1614$ & 0.65 \\
\hline
\end{tabular}
\caption{Three-model comparison over all 175 SPARC galaxies under
identical nuisance treatment ($\mu$, $D$, $i$ per galaxy with
observational priors) and the intrinsic-scatter likelihood.
$\Delta\mathrm{AIC}$ is quoted relative to the single-$a_0$ model
(negative favours the alternative); all common nuisance-parameter
counts cancel in the differences.}
\label{tab:refit}
\end{table}

The refit yields
\begin{equation}
a_0 \;=\; 0.95^{+0.03}_{-0.03} \times 10^{-10}~\mathrm{m~s^{-2}}
\quad (\Delta(-2\ln L) = 1,\ \text{profiled}),
\end{equation}
with an intrinsic velocity scatter of 5.0\% and median
$\chi^2/\mathrm{pt} = 0.83$; under floor-based error models the
fitted value moves to $1.05\times10^{-10}$ (6~km~s$^{-1}$ floor) and
$1.28\times10^{-10}$ (3~km~s$^{-1}$ floor), so the full range
$(0.96$--$1.28)\times10^{-10}$~m~s$^{-2}$ brackets the canonical
$a_0 = 1.2\times10^{-10}$.
The statistical interval is a fully profiled one: $f_{\rm int}$ is
re-optimized at each $a_0$, and the per-galaxy nuisance optimization
is polished continuously from its grid maximum (scripts
\texttt{a0\_profile\_likelihood.py}, \texttt{fitlib.py}).  The
polish matters: with grid-MAP nuisances the $-2\ln L$ profile is a
staircase with $\sim\pm10$-unit jumps---an order of magnitude above
the $\Delta = 1$ criterion---so the conditional slice quoted in an
earlier version of this analysis, $[0.945, 1.077]\times10^{-10}$,
was dominated by quantization noise, which \emph{inflated} it (the
earlier withdrawal was correct; its stated direction of bias was
not).  The smooth profile (residual jitter $0.9$ units) gives
$[0.920, 0.979]\times10^{-10}$, confirmed by an independent
curvature estimate ($[0.921, 0.978]$).
Two honesty notes accompany this interval.  First, the empirical
fold-to-fold spread of $a_0$ in the cross-validation below,
$\pm0.10\times10^{-10}$, exceeds the formal interval by a factor of
$\sim3$: the pointwise likelihood understates galaxy-level
heterogeneity (the effective sample size is closer to the 175
galaxies than to the $\sim3{,}300$ points), and we regard the fold
spread as the more conservative statistical error.  Second, the
error-model systematic above dominates both.
The stellar rescalings sit where a physical population should
($\mu$ within a factor of two of unity, no $\mu = 0$ collapse), in
sharp contrast with the free-$\mu$ Yukawa fits of
Section~\ref{sec:results}.

\textbf{Out-of-sample test.}
The single-$a_0$ claim is a claim about transfer, so we test it as
one (script \texttt{heldout\_validation.py}; 5-fold
cross-validation, folds stratified by $V_{\max}$ quartile and
quality flag).  On each fold, $a_0$ and $f_{\rm int}$ are frozen
from the training galaxies; on held-out galaxies the single-$a_0$
model may adjust only its nuisances under their observational
priors, while the Yukawa and NFW kernels refit their two per-galaxy
parameters on the test galaxy itself---an asymmetry that favours the
kernels, which are not predicting held-out data but fitting it.
Despite this handicap the single-$a_0$ model beats the Yukawa kernel
out of sample, pooled $\Delta(-2\ln L) = -4207$ (per-point $7.56$
vs $8.80$; per-galaxy wins $147/28$), and its fitted $a_0$ is stable
across folds, $(0.80$--$1.08)\times10^{-10}$ with spread
$\pm0.10$---well inside the pre-registered failure band.  NFW's
held-out advantage ($+2377$) equals its in-sample gap plus the
returned AIC penalty, as expected for a model exercising fresh
per-galaxy freedom on test data.  A split-half variant (only
$\sim86$ training galaxies) yields fold values
$[0.79, 1.42]\times10^{-10}$: the global $a_0$ estimate requires
of order a hundred galaxies to stabilize, which the 5-fold scheme
provides.  A methodological byproduct of these gates---the
grid-noise floor and its effect on interval estimation---is
documented in the repository (\texttt{results/GATES.md}).

Table~\ref{tab:refit} states the comparison.
Against the Yukawa kernel the verdict is unambiguous: one global
parameter outperforms $2\times175$ per-galaxy parameters by
$\Delta\mathrm{AIC} = 4859$, and the fixed-range kernel needs an
intrinsic scatter three times larger (15.6\%) to fit at all.
The kernel this paper set out to test is not merely disfavoured
relative to NFW; it is decisively inferior to the
acceleration-dependent response it was, per-galaxy, imitating.

Against NFW the honest statement is dual.
NFW retains a moderate total-AIC advantage,
$\Delta\mathrm{AIC} = -1614$ under the intrinsic-scatter
likelihood ($-2092$ and $-4250$ under the 6 and 3~km~s$^{-1}$
floors): the preference survives every error model, so we do not
claim parity.
Its texture, however, matters for interpretation: the per-galaxy
median is small ($\Delta\mathrm{AIC} = -2.7$; 18\% of galaxies
individually prefer the single-$a_0$ model, 25\% are
indistinguishable at $|\Delta\mathrm{AIC}| \le 2$, 57\% prefer
NFW), and roughly sixty per cent of the cumulative gap is carried
by the ten most discrepant systems---massive, high-quality rotators
where two per-galaxy parameters absorb residual systematics
(non-circular motions, spiral structure) that a zero-freedom global
law cannot.
The same numbers hold on the standard quality cut
($Q < 3$, $i \ge 30^\circ$; 153 galaxies).
A comparison against cored halo families (Einasto, DC14) under this
methodology, and a pointwise residual analysis, are the natural
next data-side steps.

\subsection{Status}

What the refit establishes within this programme is the
data-side half of the diagnosis: the SPARC sample is described,
at few-per-cent intrinsic scatter, by a single acceleration scale
consistent with $a_0 \sim \Lambda_{\rm DE}^2/M_{\rm Pl}$ up to the
factor-of-two theoretical uncertainty discussed in
Section~\ref{sec:a0}, using the very degrees of freedom
(constrained $\mu$, $D$, $i$) under which the fixed-range kernel
collapses.
What it does not establish is the theory-side half: the response
shape in Eq.~\eqref{eq:rar-kernel} is empirical, and the
programme's claim to \emph{explain} rather than reparametrize this
phenomenology rests on deriving the response from the framework's
topological sector---the pre-registered gradient-response
computation of the companion dark-energy paper, whose kill
criterion (an analytic response at $X \sim \Lambda_{\rm DE}^4$)
would close this route (Section~\ref{sec:future}).
All numbers in this section regenerate from the public repository
(scripts \texttt{rar\_refit.py}, \texttt{nuisance\_comparison.py},
\texttt{intrinsic\_scatter\_comparison.py}, each with built-in
validation gates).

\section{Conclusions}
\label{sec:conclusions}

We have confronted a Yukawa-type halo kernel, motivated by the
scalar channel of a lattice fermion theory, with 175 SPARC
galaxies.
The key findings are:
\begin{itemize}
\item In the free-$\mu$ (uninformative stellar-mass) analysis,
  58\% of galaxies yield finite cutoff radii with median
  $r_c = 15.8$~kpc and good fit quality ($\chi^2/\mathrm{pt} = 1.08$).
\item This population is prior-fragile: a physical stellar
  normalization (lognormal 0.1-dex prior, or $\mu$ fixed at the
  fiducial $\Upsilon_*$) collapses the finite-$r_c$ fraction to
  7\% and 2\% respectively, the survivors being gas-dominated
  dwarfs for which $\Upsilon_*$ is immaterial.  The free-$\mu$
  results are accordingly diagnostics of the kernel shape's
  mimicry capacity, not evidence for a screening-radius
  population.
\item The data do not support a single universal $m$: per-galaxy
  values span $\sim$1.4~decades and trend with galaxy velocity,
  pointing to an underlying \emph{acceleration scale} rather than a
  fixed length.
\item A pre-specified direct test shows, however, that the
  per-galaxy $r_c$ is not itself an acceleration-scale proxy: the
  characteristic accelerations $V^2(r_c)/r_c$ do not cluster but
  scale as $V_{\max}^{1.1}$, robustly across methodologies.  The
  single-$a_0$ structure resides at the pointwise radial
  acceleration relation, which unconstrained per-galaxy
  mass-to-light rescaling degrades; testing it within this
  framework requires the mass-to-light--constrained,
  acceleration-dependent refit performed in
  Section~\ref{sec:refit}.
\item Three independent failure modes---non-universality of $m$,
  $a_c$ tracking $V_{\max}$, and prior fragility---converge on one
  diagnosis: the Yukawa kernel is the wrong functional form, and
  its per-galaxy freedom absorbs missing acceleration-dependent
  physics.  Consistently, the companion lattice measurement places
  the scalar composite at the dynamical fermion scale
  ($m_\sigma = 2m_{\rm dyn} \gg \mathrm{kpc}^{-1}$); the last
  candidate carrier---the monopole-sector gradient response---has
  since been computed at leading order and found wanting: correct
  crossover scale, wrong exponent ($X^{1/2}$, not $X^{3/2}$), no
  galactic source.
\item This acceleration scale, $a_0 \sim \Lambda_{\rm DE}^2/M_{\rm Pl}
  \sim cH_0/2\pi$, matches the rotation-curve value within a factor of two
  and coincides with the observed dark-energy scale: an unexplained
  but robust coincidence between two \emph{observed} numbers, owing
  nothing to this framework, whose own mechanisms for both sides
  are now terminated (galactic: the leading-order response
  computations above; dark-energy: the companion paper's C2a
  measurement)---the joint-falsifiability structure of earlier
  versions, cashed in the double-negative direction.  A fixed-range
  interpretation is independently excluded by Solar-System
  fifth-force bounds.
\item Comparison with NFW shows moderate overall preference for NFW
  (unchanged in every prior variant), with reduced NFW advantage in
  dwarf galaxies in the free-$\mu$ analysis.
\item The free-$\mu$ results are stable under alternative error
  treatments, fitting-space choices, and fixed-configuration
  fitting; the stellar mass-to-light prior is the one
  methodological choice on which the finite-$r_c$ phenomenology
  decisively depends.
\item The single-$a_0$ refit (Section~\ref{sec:refit}) completes
  the data-side case: one global acceleration scale,
  $a_0 = (0.96$--$1.28)\times10^{-10}$~m~s$^{-2}$ across error
  models, with physical stellar normalizations and 5\% intrinsic
  scatter, decisively outperforms the fixed-range kernel
  ($\Delta\mathrm{AIC} = +4859$) on the very sample where the
  latter needed $2\times175$ parameters.
\item NFW retains a moderate, error-model-robust total-AIC edge
  over the single-$a_0$ model ($-1614$ to $-4250$), with a small
  per-galaxy median and the cumulative gap concentrated in a
  minority of massive systems; comparisons against cored halo
  families under identical methodology are the outstanding
  data-side test.
\end{itemize}

\subsection{Future directions}
\label{sec:future}

The central task set by this work was to determine whether the
condensate generates the acceleration-dependent (k-mouflage)
response posited in Section~\ref{sec:a0}.  Three microscopic routes
have now been executed and closed at leading order, each against a
pre-registered criterion: the pairing-channel equation of state
(analytic, $P \propto (\mu-\mu_c)^2$ with measured exponent
$2.005 \to 2.000$---a two-body theorem, never the
$(\mu-\mu_c)^{3/2}$ that superfluid dark-matter constructions
\cite{Berezhiani:2015bqa} must postulate); loop-level kinetic and
disformal couplings (Planck-suppressed,
$(b/m_{\rm dyn})^2 \sim 10^{-122}$); and the monopole-gas gradient
response (exact sine-Gordon solution: correct crossover scale,
wrong exponent, no source; Section~\ref{sec:a0}).  An
environment-triggered variant is likewise closed for every
\emph{monotone} local baryonic invariant, on five independent
grounds (magnitude, stellar consistency, local-potential degeneracy,
profile shape, ordering).  What remains registered but untested:
a \emph{non-monotonic} baryon-window response of the monopole
medium---a shell/torus-shaped activation region, geometrically
adjacent to the phantom-dark-matter structures of MOND---whose
decisive gates are a microscopically derived (not hand-drawn) window
in the monopole core action and a thick-shell response producing
flat rather than bump-like rotation curves; a near-critical
interacting monopole ensemble; and a gravitational-polarization
(anti-screening) branch.  Each is a new hypothesis, not a
correction.  Since this list was drawn up, one route of a different
type was pursued to a definite conclusion in the companion
programme: a nonlocal, topological dark-matter object whose size
would be set by the competition between the measured negative
quartic coefficient (contraction) and an \emph{endogenous}
isoscalar-vector repulsion (a Fierz-exact channel of the
microscopic interaction, dynamized with healthy kinetic term,
non-tachyonic mass, and topological-current coupling).  That
competition has now been computed, and it does \emph{not} produce a
continuum object: the vector repulsion saturates---bounded in
radius by the composite cutoff, and bounded in the coupling because
the RPA denominator screens it---while the negative-quartic
contraction diverges, so the energy collapses to a cutoff-scale
minimum with no continuum window at any coupling (companion note,
gate \texttt{P5-OMEGA-01}, \texttt{FAILED}).  This route is
therefore \emph{withdrawn as a live direction} and recorded among
the closed ones; what survives is only the distinct cutoff-scale
(lattice) question, which the effective description cannot address
and which is not a galactic-dynamics proposal.  No route on this
list currently carries microscopic support, and the paper's
conclusions rest on none of them.
On the data side, both the $r_c$-level test
(Section~\ref{sec:ac-test}) and the single-$a_0$ refit
(Section~\ref{sec:refit}) are now performed; the remaining data-side
steps are a pointwise analysis in the
$g_{\rm obs}$--$g_{\rm bar}$ plane (residual correlations of the
refit against radius, surface brightness, and gas fraction), a
comparison against cored halo profiles (Einasto, DC14) under the
same nuisance and intrinsic-scatter treatment---NFW is the hardest
two-parameter baseline in AIC terms, but not the best-fitting halo
family in the literature---the cluster regime, where MOND-type
responses historically fail and the framework must say something,
and---should any response mechanism revive---the outer-cutoff
signature that would distinguish a windowed medium from pure MOND.  The value of $\Lambda_{\rm DE}$ that fixes
$a_0$ is developed in a companion analysis of the dark-energy sector
of the same framework; closing the loop between the two would tie the
galactic acceleration scale and the cosmological constant to a single
microscopic origin.  Further extensions include a nonlinear numerical
solution of the field equation with extended baryonic sources and the
tensor channel of the same HS decomposition for emergent
gravitational dynamics.

The fitting code and data analysis pipeline are publicly available
at \url{https://github.com/zetacheng/yukawa-sparc-fits}.


% ===================================================================
\section*{Acknowledgments}
% ===================================================================

The author acknowledges the use of large language models
(Claude, Anthropic) as a research-assistance tool during this
work, including assistance with algebraic derivations and their
verification, fitting methodology, statistical analysis, and
manuscript editing.
All physical ideas and the model construction are the author's
own; all derivations and numerical results were reviewed and
verified by the author, who takes full responsibility for the
content of the manuscript.
The SPARC data were compiled by \cite{Lelli:2016zqa}.


% No appendix in this version; vacuum structure deferred to companion paper.


% ===================================================================
\begin{thebibliography}{99}

\bibitem{Lelli:2016zqa}
F.~Lelli, S.~S.~McGaugh, and J.~M.~Schombert,
``SPARC: Mass Models for 175 Disk Galaxies with Spitzer Photometry
and Accurate Rotation Curves,''
\textit{Astron.\ J.}\ \textbf{152}, 157 (2016).

\bibitem{Lelli:2017vgz}
F.~Lelli, S.~S.~McGaugh, J.~M.~Schombert, and M.~S.~Pawlowski,
``One Law to Rule Them All: The Radial Acceleration Relation of Galaxies,''
\textit{Astrophys.\ J.}\ \textbf{836}, 152 (2017).

\bibitem{McGaugh:2016leg}
S.~S.~McGaugh, F.~Lelli, and J.~M.~Schombert,
``Radial Acceleration Relation in Rotationally Supported Galaxies,''
\textit{Phys.\ Rev.\ Lett.}\ \textbf{117}, 201101 (2016).

\bibitem{Navarro:1996gj}
J.~F.~Navarro, C.~S.~Frenk, and S.~D.~M.~White,
``The Structure of Cold Dark Matter Halos,''
\textit{Astrophys.\ J.}\ \textbf{462}, 563 (1996).

\bibitem{deBlok:2009sp}
W.~J.~G.~de~Blok,
``The Core-Cusp Problem,''
\textit{Adv.\ Astron.}\ \textbf{2010}, 789293 (2010).

\bibitem{Bullock:2017xww}
J.~S.~Bullock and M.~Boylan-Kolchin,
``Small-Scale Challenges to the $\Lambda$CDM Paradigm,''
\textit{Ann.\ Rev.\ Astron.\ Astrophys.}\ \textbf{55}, 343 (2017).

\bibitem{Oman:2015xda}
K.~A.~Oman \textit{et al.},
``The unexpected diversity of dwarf galaxy rotation curves,''
\textit{Mon.\ Not.\ Roy.\ Astron.\ Soc.}\ \textbf{452}, 3650 (2015).

\bibitem{Li:2020iib}
P.~Li, F.~Lelli, S.~S.~McGaugh, and J.~M.~Schombert,
``A comprehensive catalog of dark matter halo models for SPARC galaxies,''
\textit{Astrophys.\ J.\ Suppl.}\ \textbf{247}, 31 (2020).

\bibitem{Li:2018tdo}
P.~Li, F.~Lelli, S.~S.~McGaugh and J.~M.~Schombert,
``Fitting the radial acceleration relation to individual SPARC
galaxies,''
Astron.\ Astrophys.\ \textbf{615}, A3 (2018)
[arXiv:1803.00022 [astro-ph.GA]].

\bibitem{Berezhiani:2015bqa}
L.~Berezhiani and J.~Khoury,
``Theory of dark matter superfluidity,''
Phys.\ Rev.\ D \textbf{92}, 103510 (2015)
[arXiv:1507.01019 [astro-ph.CO]].

\bibitem{Sakharov:1967pk}
A.~D.~Sakharov,
``Vacuum Quantum Fluctuations in Curved Space and the Theory of
Gravitation,''
\textit{Sov.\ Phys.\ Dokl.}\ \textbf{12}, 1040 (1968).

\bibitem{Volovik:2003fe}
G.~E.~Volovik,
\textit{The Universe in a Helium Droplet}
(Oxford University Press, 2003).

\bibitem{Barcelo:2005fc}
C.~Barcel\'{o}, S.~Liberati, and M.~Visser,
``Analogue Gravity,''
\textit{Living Rev.\ Rel.}\ \textbf{8}, 12 (2005).

\bibitem{Verlinde:2016toy}
E.~P.~Verlinde,
``Emergent Gravity and the Dark Universe,''
\textit{SciPost Phys.}\ \textbf{2}, 016 (2017).

\bibitem{Nambu:1961tp}
Y.~Nambu and G.~Jona-Lasinio,
``Dynamical Model of Elementary Particles Based on an Analogy
with Superconductivity.\ I,''
\textit{Phys.\ Rev.}\ \textbf{122}, 345 (1961).

\bibitem{Nambu:1961fr}
Y.~Nambu and G.~Jona-Lasinio,
``Dynamical Model of Elementary Particles Based on an Analogy
with Superconductivity.\ II,''
\textit{Phys.\ Rev.}\ \textbf{124}, 246 (1961).

\bibitem{Klevansky:1992qe}
S.~P.~Klevansky,
``The Nambu--Jona-Lasinio model of quantum chromodynamics,''
\textit{Rev.\ Mod.\ Phys.}\ \textbf{64}, 649 (1992).

\bibitem{Hubbard:1959ub}
J.~Hubbard,
``Calculation of Partition Functions,''
\textit{Phys.\ Rev.\ Lett.}\ \textbf{3}, 77 (1959).

\bibitem{Stratonovich:1957}
R.~L.~Stratonovich,
``On a Method of Calculating Quantum Distribution Functions,''
\textit{Sov.\ Phys.\ Dokl.}\ \textbf{2}, 416 (1957).

\bibitem{deAlmeida:2018kwq}
\'{A}.~O.~F.~de~Almeida, L.~Amendola, and V.~Niro,
``Galaxy rotation curves in modified gravity models,''
\textit{J.\ Cosmol.\ Astropart.\ Phys.}\ \textbf{2018}(08), 012 (2018).

\bibitem{Henrichs:2021bnd}
J.~Henrichs, M.~Lembo, F.~Iocco, and L.~Amendola,
``Testing gravity with the Milky Way: Yukawa potential,''
\textit{Phys.\ Rev.\ D} \textbf{104}, 043009 (2021).

\bibitem{Hassan:2025yuk}
D.~S.~Hassan, M.~D.~Danarianto, and A.~Sulaksono,
``The Milky Way and M31 rotation curves in Yukawa gravity,''
\textit{Eur.\ Phys.\ J.\ C} \textbf{86}, 198 (2026).

\bibitem{Hu:2000ke}
W.~Hu, R.~Barkana, and A.~Gruzinov,
``Fuzzy Cold Dark Matter,''
\textit{Phys.\ Rev.\ Lett.}\ \textbf{85}, 1158 (2000).

\bibitem{Hui:2016ltb}
L.~Hui, J.~P.~Ostriker, S.~Tremaine, and E.~Witten,
``Ultralight scalars as cosmological dark matter,''
\textit{Phys.\ Rev.\ D} \textbf{95}, 043541 (2017).

\bibitem{Schive:2014dra}
H.-Y.~Schive, T.~Chiueh, and T.~Broadhurst,
``Cosmic Structure as the Quantum Interference of a Coherent Dark Wave,''
\textit{Nature Phys.}\ \textbf{10}, 496 (2014).

\bibitem{Milgrom:1983ca}
M.~Milgrom,
``A modification of the Newtonian dynamics as a possible alternative
to the hidden mass hypothesis,''
\textit{Astrophys.\ J.}\ \textbf{270}, 365 (1983).

\bibitem{Famaey:2011kh}
B.~Famaey and S.~S.~McGaugh,
``Modified Newtonian Dynamics (MOND),''
\textit{Living Rev.\ Rel.}\ \textbf{15}, 10 (2012).

\bibitem{Bekenstein:2004ne}
J.~D.~Bekenstein,
``Relativistic gravitation theory for the MOND paradigm,''
\textit{Phys.\ Rev.\ D} \textbf{70}, 083509 (2004).

\end{thebibliography}

\end{document}
